% ============================================================
% Model: 一维 delta 势阱与变分法
% ============================================================

\section{一维 $\delta$ 势阱与变分法}
\label{sec:1d-delta-potential}

\begin{tipbox}[关键词标签]
宇称定理 · 基态不简并 · 无节点 · delta 势阱 · 变分法 · 试探波函数 · 归一化 · 能量期望值
\end{tipbox}

% --------------------------------------------------------
% 1. 模型定义框
% --------------------------------------------------------
\begin{modelbox}[模型：一维对称势场中的束缚态]
\textbf{物理情境}：一维势场 $V(x)=V(-x)$ 具有空间反演对称性。考察其束缚定态，特别是基态的性质。对于吸引性 $\delta$ 势阱 $V(x)=-\gamma\delta(x)$（$\gamma>0$），可用变分法近似求基态能量，并与精确解对比。

\textbf{宇称定理}：
\begin{itemize}[nosep]
    \item 若 $V(x)=V(-x)$，则哈密顿量与宇称算符对易：$[H,P]=0$
    \item 一维束缚态基态不简并，故基态波函数必为 $P$ 的本征态
    \item 基态宇称为正（偶函数），因为奇函数在原点必有节点，而基态无节点
\end{itemize}

\textbf{$\delta$ 势阱定义}：
\[
V(x)=-\gamma\delta(x),\quad\gamma>0.
\]
\begin{itemize}[nosep]
    \item 仅在 $x=0$ 处有吸引势，强度为 $\gamma$
    \item 薛定谔方程在 $x\neq0$ 处为自由粒子方程
    \item 波函数在 $x=0$ 处连续，但导数跃变：$\psi'(0^+)-\psi'(0^-)=-\frac{2m\gamma}{\hbar^2}\psi(0)$
\end{itemize}

\textbf{变分法原理}：对任意归一化试探波函数 $\psi_t$，$\langle\psi_t|H|\psi_t\rangle\geq E_0$。等号当且仅当 $\psi_t=\psi_0$（基态）时成立。
\end{modelbox}

% --------------------------------------------------------
% 2. 关键公式框
% --------------------------------------------------------
\begin{formulabox}[核心公式]
\begin{enumerate}[label=\textbf{(\arabic*)}]
    \item \textbf{宇称定理（一维基态）}：
    \[
    V(x)=V(-x)\quad\Rightarrow\quad\psi_0(-x)=\psi_0(x),\quad\psi_0(x)\neq0\;(\forall x\in\mathbb{R}).
    \]
    \texteq{基态波函数为实偶函数，无节点（除 $x\to\pm\infty$ 外）}

    \item \textbf{$\delta$ 势导数跃变条件}：
    \[
    \psi'(0^+)-\psi'(0^-)=-\frac{2m\gamma}{\hbar^2}\psi(0)
    \]
    \texteq{由薛定谔方程在 $x=0$ 邻域积分得到}

    \item \textbf{$\delta$ 势精确基态}：
    \[
    \psi_0(x)=\sqrt{\lambda}\,e^{-\lambda|x|},\quad\lambda=\frac{m\gamma}{\hbar^2},\quad
    E_0=-\frac{m\gamma^2}{2\hbar^2}.
    \]
    \texteq{仅存在一个束缚态，$E_0<0$}

    \item \textbf{变分法能量泛函}（试探波函数 $\psi_\lambda(x)=\sqrt{\lambda}\,e^{-\lambda|x|}$）：
    \[
    E(\lambda)=\langle T\rangle+\langle V\rangle=\frac{\hbar^2\lambda^2}{2m}-\gamma\lambda
    \]
    \texteq{动能项来自 $x\neq0$ 处的曲率，势能项来自 $-\gamma|\psi(0)|^2$}

    \item \textbf{变分极值}：
    \[
    \frac{dE}{d\lambda}=\frac{\hbar^2\lambda}{m}-\gamma=0\quad\Rightarrow\quad
    \lambda_{\text{opt}}=\frac{m\gamma}{\hbar^2}.
    \]
    \texteq{恰好给出精确解，说明试探波函数选取得当}

    \item \textbf{有限深势阱对比}（阱深 $V_0$，阱宽 $2a$）：
    \[
    \text{至少一个束缚态条件：}\frac{2mV_0a^2}{\hbar^2}\geq\frac{\pi^2}{8}\;(?)
    \quad\text{（一维有限深势阱\textbf{总是}至少有一个束缚态）}
    \]
    \texteq{一维势阱无论多浅多窄，总存在至少一个束缚态}
\end{enumerate}
\end{formulabox}

% --------------------------------------------------------
% 3. 训练点框
% --------------------------------------------------------
\begin{drillbox}[训练点与考法变体]
\trainingpoint{1} \textbf{宇称定理证明}：给定 $V(x)=V(-x)$，用反证法证明基态波函数必为偶函数。关键点：假设基态为奇函数 $\Rightarrow$ $\psi_0(0)=0$ $\Rightarrow$ 存在节点 $\Rightarrow$ 与基态无节点定理矛盾。

\trainingpoint{2} \textbf{变分法完整步骤}：选取试探波函数 $\to$ 归一化 $\to$ 计算 $\langle T\rangle$ 和 $\langle V\rangle$ $\to$ 构建 $E(\lambda)$ $\to$ 对变分参数求导取极值 $\to$ 代回得 $E_0$ 估计。

\trainingpoint{3} \textbf{试探波函数选取原则}：
\begin{itemize}[nosep]
    \item 满足边界条件（束缚态要求在无穷远处指数衰减）
    \item 具有正确的宇称（基态为偶函数）
    \item 含可调参数（至少一个）
    \item 形式尽量接近真实波函数（可借鉴已知类似问题的精确解）
\end{itemize}

\trainingpoint{4} \textbf{$\delta$ 势 vs 有限深势阱}：比较两者的束缚态数目、波函数形式、能量尺度。$\delta$ 势仅有一个束缚态，有限深势阱可能有多个。

\trainingpoint{5} \textbf{双 $\delta$ 势阱}：$V(x)=-\gamma[\delta(x-a)+\delta(x+a)]$，分析对称/反对称束缚态，隧穿劈裂，与分子轨道类比。
\end{drillbox}

% --------------------------------------------------------
% 4. 解题技巧框
% --------------------------------------------------------
\begin{tipbox}[快速识别与解题技巧]
\begin{itemize}
    \item \examnote{识别标志}：题干出现``一维势场''''宇称''''$\delta(x)$''''变分法''''试探波函数''''基态无节点''，立即定位本模型。
    \item \examnote{变分法步骤速查}：
    \begin{center}
    \fbox{\parbox{0.9\linewidth}{\centering
    \textbf{变分法四步曲}\\[2pt]
    1. 猜（试探波函数，含参数 $\lambda$）\\
    2. 归（归一化，确定系数）\\
    3. 算（$\langle H\rangle_\lambda=\langle T\rangle+\langle V\rangle$）\\
    4. 极（$dE/d\lambda=0$，解出 $\lambda_{\text{opt}}$）
    }}
    \end{center}
    \item \examnote{动能计算技巧}：对于分段光滑的波函数（如 $e^{-\lambda|x|}$），$d^2\psi/dx^2$ 在原点有 $\delta$ 函数奇异性。可用分部积分：
    \[
    \langle T\rangle=-\frac{\hbar^2}{2m}\int_{-\infty}^{\infty}\psi^*\frac{d^2\psi}{dx^2}dx
    =\frac{\hbar^2}{2m}\int_{-\infty}^{\infty}\Bigl|\frac{d\psi}{dx}\Bigr|^2dx
    \]
    避开二阶导数在原点的奇异性。
    \item \examnote{$\delta$ 势势能计算}：$\langle V\rangle=-\gamma\int\delta(x)|\psi(x)|^2dx=-\gamma|\psi(0)|^2$，只需波函数在原点的值。
    \item \examnote{一维 vs 三维}：一维吸引势至少有一个束缚态；三维吸引势（如球对称有限深势阱）需要足够深/宽才有束缚态。$\delta$ 势在三维需要正规化（Thomas效应）。
\end{itemize}
\end{tipbox}

% --------------------------------------------------------
% 5. 易错点框
% --------------------------------------------------------
\begin{errorbox}{常见失误与陷阱}
\begin{itemize}
    \item \textbf{动能的二阶导数陷阱}：$\psi(x)=e^{-\lambda|x|}$ 的一阶导数在 $x=0$ 处不连续：$\psi'(0^+)=-\lambda$, $\psi'(0^-)=+\lambda$。二阶导数在原点包含 $-2\lambda\delta(x)$。若直接积分 $d^2\psi/dx^2$ 而不处理奇异性，会得到错误结果。
    \item \textbf{基态简并的维度依赖}：一维束缚态基态不简并；二维及以上基态可能简并（如氢原子基态 $l=0,m=0$ 不简并，但激发态可简并）。宇称定理的``基态必为偶函数''仅在一维成立。
    \item \textbf{试探波函数必须归一化}：变分原理要求试探波函数归一化。若忘记归一化常数，能量期望值会偏差。
    \item \textbf{变分法给出的是上界}：$E(\lambda_{\text{opt}})\geq E_0$。若试探波函数选得好，接近精确解；选得差，上界很松。
    \item \textbf{$\delta$ 势的维度}：一维 $V(x)=-\gamma\delta(x)$ 有且仅有一个束缚态；二维 $V(\mathbf{r})=-\gamma\delta^{(2)}(\mathbf{r})$ 有一个束缚态（对数发散需正规化）；三维 $V(\mathbf{r})=-\gamma\delta^{(3)}(\mathbf{r})$ 无束缚态（Thomas崩溃）。
    \item \textbf{宇称定理的反用}：若已知基态波函数有节点，则可反推势场不对称或存在外场破坏空间反演对称性。
\end{itemize}
\end{errorbox}

% --------------------------------------------------------
% 6. 物理图像与关联模型
% --------------------------------------------------------
\begin{insightbox}[物理图像与模型关联]
\textbf{物理图像}：

$\delta$ 势阱就像一个无限窄、无限深的陷阱：
\begin{itemize}[nosep]
    \item 粒子在 $x\neq0$ 处完全自由，像``在旷野中游荡''。
    \item 一旦靠近 $x=0$，就被瞬间吸入陷阱——但只有足够``轻''（动量小）的粒子才能被束缚。
    \item 波函数呈指数衰减 $e^{-\lambda|x|}$，衰减长度 $1/\lambda=\hbar^2/m\gamma$。势越强（$\gamma$ 越大），粒子越局域。
    \item 整个系统只有一个束缚态——无论势多强，深能级只有一个。
\end{itemize}

\textbf{基态无节点定理的证明思路}：

若基态 $\psi_0(x)$ 在某处 $x=a$ 有节点（$\psi_0(a)=0$），可构造新波函数 $\phi(x)=|\psi_0(x)|$。
\begin{itemize}[nosep]
    \item $\phi$ 与 $\psi_0$ 概率密度相同，故归一化相同
    \item 动能项 $\langle T\rangle_\phi\leq\langle T\rangle_{\psi_0}$（因为 $|\psi|$ 的``尖点''使动能更小，或通过分部积分证明）
    \item 势能项相同（$\langle V\rangle$ 只依赖于 $|\psi|^2$）
    \item 故 $\langle H\rangle_\phi\leq\langle H\rangle_{\psi_0}=E_0$，但 $E_0$ 是最低能量，矛盾
    \item 因此基态波函数不能有节点
\end{itemize}

\textbf{与相似模型的对比}：
\begin{center}
\begin{tabular}{llll}
\toprule
\textbf{对比维度} & \textbf{$\delta$ 势阱} & \textbf{有限深势阱} & \textbf{谐振子} \\
\midrule
束缚态数目 & 1 个 & $N$ 个（取决于 $V_0a^2$） & 无限多个 \\
基态波函数 & 指数衰减 $e^{-\lambda|x|}$ & 余弦（阱内）+ 指数（阱外） & 高斯 $e^{-m\omega x^2/2\hbar}$ \\
能谱 & 仅离散（1个） & 离散 + 连续 & 等间距 $E_n=\hbar\omega(n+1/2)$ \\
变分法效果 & 试探波函数恰为精确解 & 给出上界估计 & 高斯试探波函数给出精确基态 \\
\bottomrule
\end{tabular}
\end{center}

\textbf{推广方向}：
\begin{itemize}[nosep]
    \item 双 $\delta$ 势阱：对称/反对称束缚态，能级隧穿劈裂，与氢分子离子 $H_2^+$ 类比
    \item 周期 $\delta$ 势（Kronig-Penney 模型极限）：$V(x)=-\gamma\sum_n\delta(x-na)$，能带结构
    \item 三维 Thomas 效应：吸引 $\delta$ 势在三维导致能量无下界，需引入紫外截断
    \item 非线性薛定谔方程中的 $\delta$ 势：孤子与束缚态的相互作用
\end{itemize}
\end{insightbox}

% --------------------------------------------------------
% 7. 推导附录
% --------------------------------------------------------
\begin{derivationbox}[推导细节（自我检验用）]
\textbf{Step 1：宇称定理严格证明}

已知 $[H,P]=0$，$H$ 与 $P$ 有共同本征函数完备集。设 $\psi_0(x)$ 为基态波函数，$E_0$ 为基态能量。

假设 $\psi_0$ 为奇函数：$\psi_0(-x)=-\psi_0(x)$。则 $\psi_0(0)=-\psi_0(0)=0$。

构造 $\phi(x)=|\psi_0(x)|$。$\phi$ 为偶函数，无节点（除无穷远外），且归一化与 $\psi_0$ 相同。

计算 $\langle H\rangle_\phi$：
\[
\langle H\rangle_\phi=\langle T\rangle_\phi+\langle V\rangle_\phi.
\]

势能：$\langle V\rangle_\phi=\int V(x)|\phi|^2dx=\int V(x)|\psi_0|^2dx=\langle V\rangle_{\psi_0}$（因 $V$ 只依赖于概率密度）。

动能（分部积分）：
\[
\langle T\rangle_\phi=\frac{\hbar^2}{2m}\int\Bigl|\frac{d\phi}{dx}\Bigr|^2dx.
\]

在 $\psi_0$ 的节点之间，$\phi=\pm\psi_0$，$|d\phi/dx|=|d\psi_0/dx|$。但在节点处，$\psi_0$ 穿过零点，$\phi$ 形成尖点，$d\phi/dx$ 在节点处不连续。对于基态（只有一个节点在原点），可分段证明 $\langle T\rangle_\phi\leq\langle T\rangle_{\psi_0}$。

因此 $\langle H\rangle_\phi\leq E_0$。但 $E_0$ 是最低能量本征值，矛盾。故 $\psi_0$ 不能是奇函数，必为偶函数。

\textbf{Step 2：$\delta$ 势精确解}

薛定谔方程（$x\neq0$）：
\[
-\frac{\hbar^2}{2m}\frac{d^2\psi}{dx^2}=E\psi,\quad E<0.
\]

令 $\kappa=\sqrt{-2mE}/\hbar$，解为 $\psi(x)=Ae^{-\kappa|x|}$（束缚态要求指数衰减，且基态为偶函数）。

归一化：
\[
\int_{-\infty}^{\infty}|\psi|^2dx=2|A|^2\int_0^{\infty}e^{-2\kappa x}dx=\frac{|A|^2}{\kappa}=1
\quad\Rightarrow\quad A=\sqrt{\kappa}.
\]

导数跃变条件：
\[
\psi'(0^+)-\psi'(0^-)=-2\kappa\sqrt{\kappa}=-\frac{2m\gamma}{\hbar^2}\sqrt{\kappa}.
\]

故
\[
\kappa=\frac{m\gamma}{\hbar^2},\quad E_0=-\frac{\hbar^2\kappa^2}{2m}=-\frac{m\gamma^2}{2\hbar^2}.
\]

\textbf{Step 3：变分法计算}

试探波函数：$\psi_\lambda(x)=\sqrt{\lambda}\,e^{-\lambda|x|}$（已归一化）。

动能（分部积分法）：
\[
\langle T\rangle=\frac{\hbar^2}{2m}\int_{-\infty}^{\infty}\Bigl|\frac{d\psi_\lambda}{dx}\Bigr|^2dx
=\frac{\hbar^2}{2m}\cdot2\lambda^2\cdot\lambda\int_0^{\infty}e^{-2\lambda x}dx
=\frac{\hbar^2\lambda^2}{2m}.
\]

势能：
\[
\langle V\rangle=-\gamma|\psi_\lambda(0)|^2=-\gamma\lambda.
\]

总能量：
\[
E(\lambda)=\frac{\hbar^2\lambda^2}{2m}-\gamma\lambda.
\]

求导：
\[
\frac{dE}{d\lambda}=\frac{\hbar^2\lambda}{m}-\gamma=0\quad\Rightarrow\quad\lambda_{\text{opt}}=\frac{m\gamma}{\hbar^2}=\kappa.
\]

代回：
\[
E(\lambda_{\text{opt}})=\frac{\hbar^2}{2m}\Bigl(\frac{m\gamma}{\hbar^2}\Bigr)^2-\gamma\Bigl(\frac{m\gamma}{\hbar^2}\Bigr)
=-\frac{m\gamma^2}{2\hbar^2}=E_0.
\]

变分法恰好给出精确解，说明试探波函数的函数形式与真实基态完全一致。
\end{derivationbox}

\vspace{1em}
